\chapter{Conclusion}
\label{chap:conclusion}

On-board processing exists to shorten the path from acquisition to actionable information for time-critical Earth observation applications such as flood and wildfire response, where the value of a detection decays with every hour spent waiting for a ground-based product.
Because a mission's own sensor data cannot exist before launch, and labeled examples of rare events accumulate slowly afterward, foundation models pre-trained on other missions' archives are an attractive way to reach useful performance with few mission-specific labels.
This attractiveness is conditional, however, on a question this thesis set out to answer directly: whether such a model's learned representation actually survives contact with a sensor it was never trained on, and, if not, whether that can be fixed deliberately instead of assumed away.

The evidence gathered here says it can, but not for free, and not by the mechanism this thesis originally expected.
TerraMind's latent representation of real $\Phi$-sat-2 imagery does differ measurably from its representation of Sentinel-2 imagery of the same scenes: the two sensor views of the same ground are close to orthogonal in its embedding space, and a linear probe transfers only a third of its source-domain accuracy to the target sensor. This separation is present from the encoder's very first layer, pointing to a low-level, radiometric origin, consistent with differing spectral response between the two sensors, and not to a semantic effect that would be expected to emerge only in deeper layers.
Naively fine-tuning on physics-simulated degraded data does not close this gap; an explicit domain-adaptation mechanism is required, and the co-registered sensor pairs collected for this thesis made it possible to test candidate mechanisms against real data instead of assuming one would work. The mechanism that did close the gap, domain-specific batch normalization, needs only to know which sensor produced a given input, not that inputs be paired across sensors; a paired-contrastive objective that exploits the pairing directly was also implemented and added nothing measurable beyond it on the LULC task. Both mechanisms were applied during knowledge distillation, the same step that is separately required to compress the model onto space-qualified hardware that cannot run its original Vision Transformer architecture at all, so that the deployable student is domain-adaptive by construction, not by hoped-for inheritance from the teacher.
Whether this adaptation is actually worth its cost is answered conditionally rather than unconditionally: distilling with domain-specific normalization outperforms a supervised baseline trained directly on the target sensor at a representative label budget, but loses to that same baseline at a much smaller one, reversing in both directions, though the small-budget end of that reversal falls within the run-to-run variance this thesis is able to resolve. Treating the domain-gap problem and the compute-constraint problem as one combined training objective, and not as two sequential and independent fixes, is the thesis's central methodological position, with the important qualification that whether this combination is worth adopting depends on how many target-domain labels are actually available.

This work is limited to a single mission and sensor pair, and to one specific pre-trained GeoFM; whether the spectral-response-driven account of the domain gap generalizes to other on-board missions, and whether TerraMind's broad pre-training is in fact more label-efficient than a mission-specific model trained from scratch, are both open questions this thesis frames precisely but does not close.
The comparison against $\Phi$satNet, a bespoke GeoFM developed for the same mission by the same research group, is the immediate next step toward closing the second of these, and a natural extension of this work would repeat the analysis for other sensor pairs and other missions, or attempt to train a custom foundation model instead of adapting an existing one, to test whether the findings here are specific to TerraMind or general to the class of large, diversely pre-trained EO models it represents.
